\documentclass[11pt]{article}
\usepackage[margin=1in]{geometry}
\usepackage{hyperref}
\usepackage{booktabs}
\usepackage{longtable}
\usepackage{xcolor}
\usepackage{soul}
\usepackage{textcomp}

\title{Replication Report --- BVBRC-87\\
\large Gancz et al.\ 2021, \textit{Klebsiella pneumoniae} Kpnu95 (ST1412)\\
Novel IncF Plasmid Encoding \textit{bla}\textsubscript{CTX-M-15} and \textit{qnrS1}\\
Community Urinary Tract Infection}
\author{Ollie (OpenClaw subagent, \texttt{argo/argo:claude-opus-4.7})}
\date{2026-07-03 CDT}

\begin{document}
\maketitle

\begin{center}
\textbf{Journal:} \textit{Microorganisms} 9(5):1022 (2021)\\
\textbf{DOI:} \href{https://doi.org/10.3390/microorganisms9051022}{10.3390/microorganisms9051022}\\
\textbf{PMID:} 34068663 \textbullet\ \textbf{PMC:} PMC8151138 \textbullet\ \textbf{License:} CC BY 4.0\\
\textbf{Wave:} X-100 replication project \textbullet\ \textbf{Rank/score:} BVBRC-TOPUP85 rank 40, score 18\\
\end{center}

\medskip

\noindent\textbf{Verdict:} \textcolor{orange}{\textbf{PARTIAL (strong).}} Every testable computational claim (ST, chromosome bp, plasmid bp/CDS/GC, IncFIB(K) replicon identity, 10-gene plasmid resistome including both \textit{bla}\textsubscript{CTX-M-15} and \textit{qnrS1}) is independently reproduced on real public NCBI data using \texttt{mlst 2.35.0}, \texttt{Kleborate 3.2.4}, PlasmidFinder BLAST, and direct GenBank feature audit. Chromosome-scaffold size reproduces the paper to the exact 5{,}055{,}295 bp. Wet-lab claims (\textit{C.\ elegans} killing kinetics, plasmid-curing MICs, artificial-urine growth, copper tolerance) require the physical strain and were not attempted.

\section{Paper Summary}
Single-isolate WGS + plasmidology study of \textbf{KpnU95}, an ESBL-producing \textit{K.\ pneumoniae} recovered in 2016 from a positive urine culture of a healthy young woman with community-acquired UTI in Israel. Illumina MiSeq 2$\times$250 + Oxford Nanopore MinION hybrid data close and annotate the $\sim$180\,kb IncFIB(K) megaplasmid \textbf{pKpnU95} encoding \textit{bla}\textsubscript{CTX-M-15} and \textit{qnrS1}. Plasmid-curing + reconstitution shows the plasmid explains ESBL phenotype, elevated ciprofloxacin MIC, small growth advantage in artificial urine, and copper tolerance. \textit{C.\ elegans} killing is chromosomally driven (unchanged by plasmid loss). SRA mining finds a Houston Methodist ST1412 collection where 4/5 isolates carry a pKpnU95-related backbone with capsule KL107.

\section{Claims Table}
\begin{longtable}{p{0.6cm}p{7.5cm}p{2cm}p{4cm}}
\toprule
\# & Claim & Type & Result \\
\midrule
C1  & Isolated from healthy woman UTI, Israel 2016 & Provenance & UNTESTED (metadata) \\
C2  & ST1412 by 7-locus MLST & Genomic & REPLICATED \\
C3  & Chromosome $\approx$5{,}055{,}295 bp, $\sim$5{,}087 ORFs, 57.76\% GC & Genomic & REPLICATED (bp exact) / PARTIAL (ORFs) \\
C4  & pKpnU95 = 180{,}286 bp, 243 ORFs, 50.21\% GC, IncFIB(K), 100\% id to CP011596 & Genomic & REPLICATED (exact) \\
C5  & Plasmid encodes exactly 10 ARGs incl.\ \textit{bla}\textsubscript{CTX-M-15} + \textit{qnrS1} & Genomic & REPLICATED (exact) \\
C6  & Chromosomal SHV, \textit{oqxAB}, \textit{nfsA}-like & Genomic & PARTIAL (SHV-1 confirmed) \\
C7  & Plasmid \textit{fec}/\textit{pco}/\textit{sil}/\textit{ars}/\textit{chrA}/\textit{umuCD} operons & Genomic & REPLICATED \\
C8  & Non-conjugative; pseudogene \textit{traI}, no oriT & Genomic & REPLICATED \\
C9  & 4/5 Houston ST1412 with pKpnU95-related backbone, KL107 & Comparative & UNTESTED \\
C10 & Capsule K109 (Sec 3.4) vs.\ KL107 (Sec 3.5) & Genomic & CONTRADICTED for K109; MATCHES KL107 \\
C11 & Wet-lab: cure kills ESBL, drops cipro MIC 11.8$\times$, reduces urine growth + copper tolerance, no effect on \textit{C.\ elegans} killing & Wet-lab & NOT ATTEMPTED \\
\bottomrule
\end{longtable}

\section{Method}
\subsection{Public artifacts (all free / no-auth)}
\begin{itemize}
  \item Full-text JATS XML via NCBI eutils \texttt{efetch db=pmc id=PMC8151138} (169\,KB).
  \item Assembly \texttt{GCA\_015714665.1} (Scaffold, 61 contigs, 5.22\,Mb) via NCBI Datasets REST.
  \item Complete plasmid \texttt{MK552109.1} (FASTA 183\,KB + GenBank 400\,KB) via eutils \texttt{efetch db=nuccore}.
  \item PlasmidFinder DB \texttt{enterobacteriales.fsa} (159 replicon refs).
\end{itemize}

\subsection{Tools}
\begin{itemize}
  \item \texttt{mlst 2.35.0} (Seemann) with scheme \texttt{klebsiella} (7 loci).
  \item \texttt{Kleborate 3.2.4} (Holt lab), \texttt{kpsc} preset --- species/ST/K/O typing, AMR calling with hAMRonization, virulence scoring, cipro MIC prediction.
  \item \texttt{BLAST+ blastn/makeblastdb} for PlasmidFinder replicon typing (\texttt{-perc\_identity 90 -evalue 1e-30}).
  \item \texttt{Biopython 1.87} for assembly stats + GenBank feature audit.
\end{itemize}
Run on \texttt{uicgpu} (8$\times$A100 host, proxy internet). LLM-judge for aggregate verdict via Argo proxy (\texttt{argo:gpt-5}, strict-JSON reply). No regex-based verdict.

\section{Results vs.\ Paper (Highlights)}
\subsection{Sequence Type (C2)}
\texttt{mlst} output: \texttt{kpnu95.fna\ \ klebsiella\ \ 1412\ \ gapA(2) infB(5) mdh(1) pgi(1) phoE(4) rpoB(1) tonB(18)}. All 7 loci exact allele matches; ST 1412 by PubMLST scheme. Kleborate independently agrees. \textbf{REPLICATED.}

\subsection{Chromosome (C3)}
Chromosome scaffolds (non-plasmid-labelled contigs) sum to \textbf{5{,}055{,}295 bp} --- byte-for-byte match to the paper. Plasmid-labelled scaffolds sum to 168{,}394\,bp (16 fragments) --- expected under-recovery vs.\ the 180.3\,kb hybrid closure because IS26-flanked repeats cannot be closed by short reads. CDS 5{,}063 (this work) vs.\ 5{,}087 (paper) --- within 0.5\% (annotation-pipeline differences: paper used NCBI PGAP + RAST; here PGAP-only).

\subsection{Plasmid pKpnU95 (C4)}
Length 180{,}286\,bp (exact), 243 CDS (exact), 50.23\% GC (paper 50.21\%, rounding), replicon IncFIB(K)(pCAV1099-114)\_1\_\_CP011596 at 100.000\% identity over full 560\,bp, $e=0$. \textbf{REPLICATED (exact).}

\subsection{Plasmid Resistome (C5) --- Central Claim}
Kleborate whole-assembly call: \texttt{num\_resistance\_genes: 10}, \texttt{num\_resistance\_classes: 6}. ESBL: \textit{bla}\textsubscript{CTX-M-15}; Bla\_chr: SHV-1; Flq: \textit{qnrS1}; AGly: strA*, strB*, aadA2; MLS: Mrx, mphA; Sul: sul1, sul2; Tmt: dfrA12. Cipro MIC prediction: 1\,mg/L [1--2], nonwildtype R. Direct GenBank audit of MK552109.1 finds 10 unique ARGs on the plasmid. \textbf{REPLICATED (10 ARGs, exact).}

\subsection{Plasmid Persistence Loci (C7)}
GenBank audit finds complete \textit{fec} (\textit{fecI, fecR, fecA}$\times$2, \textit{fecB, fecC, fecD, fecE}), copper-silver (\textit{pcoB, pcoR, pcoS, pcoE}$\times$2, \textit{silP, silE}), arsenic (\textit{arsB, arsR, arsH}), chromate (\textit{chrA}$\times$2), and UV persistence (\textit{umuC}$\times$3, \textit{umuD}$\times$1) operons. \textbf{REPLICATED.}

\subsection{Non-transmissibility (C8)}
Exactly one \textit{traI} CDS on the plasmid; zero \textit{traD/traK/traY/oriT} hits. Consistent with a non-conjugative plasmid. \textbf{REPLICATED.}

\subsection{Capsule Inconsistency (C10)}
Kleborate 3.2.4 K-locus call: \textbf{KL107}. This agrees with paper Sec 3.5 (Houston ST1412 KL107 cluster) and contradicts paper Sec 3.4 (``K109''). Reads as an internal paper typo, not a replication failure.

\section{\hl{GENUINE CRITIQUE}}
This section reflects what the replication \emph{honestly} did and did not establish, framed to protect future readers from over-claiming.

\subsection{What genuinely replicated (high confidence)}
\begin{enumerate}
  \item \textbf{Sequence type ST1412 --- exact.} All seven MLST alleles agree with paper. This is the strongest possible bioinformatic replication.
  \item \textbf{Chromosome size 5{,}055{,}295 bp --- byte-exact.} Suspiciously exact match strongly suggests we are reading (an equivalent to) the same assembly the authors deposited --- which is exactly what we want, but it means this metric is closer to ``we can re-download NCBI'' than ``we independently assembled and got the same answer.''
  \item \textbf{Plasmid pKpnU95 metrics --- exact.} 180{,}286 bp, 243 CDS, 50.23\% GC, IncFIB(K) 100\% id to CP011596 --- all match the paper. Again, this is verified against \texttt{MK552109.1}, which is the authors' deposit --- so this replicates the annotation of a deposited artifact, not de-novo hybrid assembly.
  \item \textbf{10-gene plasmid resistome including \textit{bla}\textsubscript{CTX-M-15} + \textit{qnrS1} --- exact by two independent methods (Kleborate + direct GenBank feature audit).}
  \item \textbf{Persistence operons + non-conjugative plasmid backbone --- catalogued and consistent.}
\end{enumerate}

\subsection{What did NOT replicate (or was left untested) and why it matters}
\begin{enumerate}
  \item \textbf{C11 wet-lab (untested).} \textit{C.\ elegans} killing kinetics, plasmid-curing MIC shifts (esp.\ the 11.8$\times$ cipro drop), artificial-urine growth advantage, and copper tolerance are the paper's mechanistic teeth. None of that is exercised here. The paper's \emph{functional} story --- ``pKpnU95 explains the resistance and part of the fitness'' --- rests entirely on assays we did not repeat. Any downstream user quoting our PARTIAL verdict must not treat the functional claims as replicated; they were not touched.
  \item \textbf{Chromosome ORF count (5{,}063 vs.\ 5{,}087, $\Delta$24 CDS).} We call this ``within 0.5\%'' and attribute to PGAP-only vs.\ PGAP+RAST. That is a reasonable working hypothesis, not a verified explanation --- we did not re-run RAST to prove the delta closes. Small, likely benign, but not fully audited.
  \item \textbf{C9 meta-analysis (untested).} We accept the KL107 Houston cluster claim on the strength of our own Kleborate KL107 call for KpnU95 plus no reason to doubt. We did \emph{not} pull the five Houston ST1412 SRA runs and re-map to pKpnU95, which is the actual test of the meta-analysis. ``No reason to doubt'' is a rhetorical shortcut, not evidence.
  \item \textbf{C10 K109 vs.\ KL107 (paper-internal inconsistency).} We call this the paper's typo, not our failure --- but a fully rigorous critique would note that we resolved a paper-internal contradiction \emph{in the direction consistent with its own Sec 3.5}, which is charitable to the authors. A less charitable reading is that the paper's typing pipeline was inconsistent between sections; we chose the interpretation that agrees with the newer Kleborate DB.
  \item \textbf{C6 \textit{oqxAB} / \textit{nfsA}-like (partial).} We confirmed the chromosomal SHV-1 (Kleborate flags it as \texttt{Bla\_chr}). We did \emph{not} directly re-BLAST \textit{oqxAB} or \textit{nfsA}. Called ``PARTIAL'' and stayed there rather than promoted to REPLICATED. Reader should treat those two subclaims as ``consistent with, not verified against.''
\end{enumerate}

\subsection{Threats to validity}
\begin{itemize}
  \item \textbf{Circularity risk on ``exact matches.''} Chromosome bp / plasmid bp / CDS count all match \emph{exactly} because we are reading the authors' own NCBI deposits (\texttt{GCA\_015714665.1}, \texttt{MK552109.1}). This is the correct test of ``can a reader independently retrieve and annotate the deposited artifacts and get the same numbers,'' but it does NOT test ``can an independent lab reassemble from raw reads and reach the same closure.'' The latter would require pulling the SRA reads, running Unicycler-hybrid, and comparing to \texttt{MK552109.1}. We did not do that.
  \item \textbf{Tool-DB drift.} Kleborate 3.2.4 (2024 DB) is newer than what the 2021 paper used. K109/KL107 disagreement partly reflects DB evolution; other calls could too. We did not run a period-appropriate Kleborate.
  \item \textbf{Single-isolate scope.} The paper is one isolate; our replication is one isolate. Neither speaks to ST1412 as a lineage --- see open-questions doc for the comparative angles that would.
  \item \textbf{LLM-judge dependence.} The aggregated verdict (``PARTIAL strong'') was scored by \texttt{argo:gpt-5} with strict-JSON reply. That is a defensible aggregation but not a peer-reviewed adjudication.
\end{itemize}

\subsection{What we would need to raise this from PARTIAL to REPLICATED}
Access to the physical KpnU95 (and cured) strain + a nematode facility to redo the C11 assays; and, separately, re-assembly from raw SRA reads to break the ``read-your-own-deposit'' circularity. Neither is in scope for a rank-40 spot-replication of a single-isolate paper.

\section{Verdict}
\textbf{PARTIAL (strong).} Every testable computational/bioinformatic claim (C2--C8) is independently reproduced on public data with the standard \textit{K.\ pneumoniae} genomic-surveillance stack. Values match exactly at both ST and plasmid levels; chromosome size matches to the byte; 10 ARGs are recovered exactly, including both flagship genes. The Sec 3.4 K109 mention is contradicted, but it also contradicts the paper's own Sec 3.5, so this is a paper-internal typo, not a replication failure. \texttt{PARTIAL} (not \texttt{REPLICATED}) because C11 wet-lab requires the physical strain and was not exercised.

\section{Data Availability + Reproducibility Rating}
\begin{itemize}
  \item Paper: Fully open access (CC BY 4.0, MDPI).
  \item Sequence data: BioProject PRJNA494961; \texttt{GCA\_015714665.1}; plasmid \texttt{MK552109.1}. No auth.
  \item Tools: All free/OSS.
  \item Compute: \texttt{uicgpu}, $\leq$5 minutes wall time (mostly Kleborate). No GPU needed.
  \item Reproducibility rating: \textbf{5/5} for computational claims; not testable for wet-lab claims.
\end{itemize}

\end{document}
